\begin{center}
\Huge
	Komplekse andengradspolynomier
\end{center}
\section*{Polynomier med komplekse løsninger}
\stepcounter{section}

Vores motivation var fra starten at løse ligninger på formen $z^2 + 1 = 0$, og vi indførte derfor vores komplekse enhed $i = \sqrt{-1}$. Det viser sig, at alle andengradspolynomier har to komplekse løsninger talt med multiplicitet. Dette er et specialtilfælde af det, vi kalder for algebraens fundamentalsætning
\begin{setn}[Algebraens fundamentalsætning]
	Ethvert polynomium $p$ givet ved
	\begin{align*}
		p(z) = a_nz^n + a_{n-1}z^{n-1} + \cdots + a_1 z + a_0,
	\end{align*}
	hvor $a_1, a_2, \hdots, a_n \in \mathbb{C}$ og $a_n \neq 0$ har $n$ komplekse rødder talt med multiplicitet. 
\end{setn}
Det er ganske besværligt at bevise denne sætning, og man kan sagtens tilrettelægge et helt kursus på universitetet, der som mål skal vise sætningen. Vi nøjes derfor blot med at påstå sætningen. Vi vil heller ikke definere præcist hvad multipliciteten af en rod betyder, da dette er lettest at beskrive, hvis man har set polynomisk division. 

Vi vil finde rødder til komplekse andengradspolynomier nøjagtigt som vi gjorde det med reelle andengradspolynomier. Vi opskriver dette som en sætning først.
\begin{setn}[Rodformlen for komplekse andengradspolynomier]\label{setn:rodformel}
	Et komplekst andengradspolynomium $p$ på formen
	\begin{align*}
		p(z) = az^2 + bz + c
	\end{align*}
	hvor $a,b,c\in \mathbb{C}$ og $a \neq 0$ har de komplekse rødder
	\begin{align*}
		z_1 &= \frac{-b + \sqrt{b^2 - 4ac}}{2a},\\
		z_2 &= \frac{-b - \sqrt{b^2 - 4ac}}{2a}.\\
	\end{align*}
\end{setn}
\begin{exa}
	Vi betragter polynomiet $p$ givet ved
	\begin{align*}
		p(z) = z^2 + 2z + 2.
	\end{align*}
	Vi bestemmer rødderne ved at bruge diskriminantformlen som altid. Først bestemmes diskriminanten
	\begin{align*}
		d = b^2 - 4ac = 2^2 - 4\cdot 1\cdot 2 = 4 - 8 = -4.
	\end{align*}
	Da diskriminanten er negativ, har vi ingen reelle rødder, men dette stopper os ikke mere! Vi sætter ind i formlen 
	for rødderne
	\begin{align*}
		z_1 &= \frac{-b + \sqrt{d}}{2a} = \frac{-2 + \sqrt{-4}}{2} = \frac{-2+\sqrt{4}\sqrt{-1}}{2} 
		= \frac{-2 + 2i}{2} = -1+i\\
		z_2 &=\frac{-b - \sqrt{d}}{2a} = \frac{-2 - 2i}{2} = -1 - i. 
	\end{align*}
	De komplekse rødder til $p$ er derfor $z_1 = -1 + i $ og $z_2 = -1 - i$. 
\end{exa}

\subsection*{Opgave 1}
Bestem rødderne for følgende komplekse andengradspolynomier
\begin{align*}
	&1) \ p(z) = z^2 - 4z + 5 &&2) \ p(z) = z^2 + 2z + 10 \\
	&3) \ p(z) = -z^2 + 6z - 13 &&4) \ p(z) = z^2 + 4z + 8 \\
	&5) \ p(z) = -z^2 - 2z - 17 &&6) \ p(z) = z^2 - 8z + 25
\end{align*}

\subsection*{Opgave 2}
Af Opgave 1 kan vi se, at hvis koefficienterne for $p$ er reelle, så vil rødderne for $p$ være parvist komplekskonjugerede. Vis, hvorfor dette er tilfældet ud fra Sætning \ref{setn:rodformel}.

\subsection*{Opgave 3}
Vis, at hvis et polynomium er $p$ givet ved
\begin{align*}
	p(z) = z^2 + bz + c
\end{align*}
hvor $b$ og $c$ begge er reelle, så opfylder rødderne $z_1 = \alpha + \beta i$ og $z_2 = \alpha - \beta i$, at 
\begin{align*}
	b = 2\alpha 
\end{align*}
og 
\begin{align*}
	c = \alpha^2 + \beta^2
\end{align*}
(Vink: Opskriv den faktorisede form af $p$.)

\subsection*{Opgave 4}
Bestem rødderne for følgende polynomier med komplekse koefficienter. 
\begin{align*}
	&1) \  p(z) = z^2 + iz + i -1   &&2) \ p(z) = 2z^2 + (6-4i)z - 12i  \\
	&3) \  p(z) = z^2 +(-2+5i)z -2i - 4   &&4) \ p(z) = z^2-2iz-1+8i \\ 
\end{align*}

\subsection*{Opgave 5}
Vi vil gerne vise, at hvis $z_0$ er en rod til $p$ givet ved
\begin{align*}
	p(z) = z^n + a_{n-1}z^{n-1} + \cdots + a_1z + a_0,
\end{align*}
hvor koefficienterne er reelle, så er $\overline{z_0}$ også en rod til $p$.
\begin{enumerate}[label=\roman*)]
	\item Vis, at $\overline{z_1} + \overline{z_2} = \overline{z_1 + z_2}$ for to komplekse tal $z_1$ og $z_2$.
	\item Vis, at $\overline{z_1}\overline{z_2} = \overline{z_1z_2}$.
	\item Brug dette til at argumentere for, at $\overline{z}^n = \overline{z^n}$ for alle $n \in \mathbb{N}$. 
	\item Vis, at $r\overline{z_1} = \overline{rz_1}$ for reelle tal $r$. 
	\item Brug nu dette til at vise, at $\overline{z}$ også er en rod til $p$. 
\end{enumerate}
\subsection*{Opgave 6}
Brug resultatet fra Opgave 5 samt algebraens fundamentalsætning til at argumentere for, at polynomier med reelle koefficienter af ulige grad må have mindst én reel rod. 
